\chapter{Conclusiones y desarrollos futuros}\label{conclusiones}

\section{S\'intesis final del trabajo}

El desarrollo realizado ha permitido construir un sistema web de gesti\'on de
entregables acad\'emicos con una orientaci\'on claramente profesional,
integrando an\'alisis formal de requisitos, dise\~no de arquitectura,
implementaci\'on por capas, verificaci\'on en profundidad y despliegue continuo.

El resultado no se limita a un prototipo funcional,
sino a una soluci\'on trazable y mantenible,
alineada con el problema de negocio detectado en el contexto docente:
mejorar la gesti\'on de actividades, entregas, evaluaciones y feedback
frente a herramientas generalistas con menor adecuaci\'on a la operativa diaria.

\section{Grado de cumplimiento de objetivos}

El objetivo general planteado en la memoria
(dise\~nar e implementar una plataforma de gesti\'on de entregables usable,
segura y evolutiva) puede considerarse \textbf{cumplido}.

Respecto a los objetivos espec\'ificos, el balance es el siguiente:

\begin{itemize}
	\item \textbf{Modelado del dominio acad\'emico}: cumplido.
	Se han implementado entidades, relaciones y reglas de negocio acordes al ERS/DAS.
	\item \textbf{Soporte completo del flujo de entrega/evaluaci\'on}: cumplido.
	El sistema cubre ciclo de vida de entregables, entregas, calificaci\'on y feedback.
	\item \textbf{Control de acceso por roles y contexto}: cumplido.
	Se ha materializado autorizaci\'on con verificaciones por perfil y recurso objetivo.
	\item \textbf{Calidad t\'ecnica y verificaci\'on}: cumplido.
	Se dispone de estrategia de pruebas en capas con m\'etricas elevadas de cobertura y mutaci\'on.
	\item \textbf{Despliegue y operaci\'on automatizada}: cumplido.
	Se han definido pipelines CI/CD, controles de seguridad y configuraci\'on por entornos.
\end{itemize}

\section{Aportaciones principales}

Las aportaciones m\'as relevantes del trabajo son:

\begin{enumerate}
	\item \textbf{Aportaci\'on funcional}:
	una plataforma centrada en la operativa docente,
	con flujos diferenciados para administraci\'on, profesorado y alumnado.
	\item \textbf{Aportaci\'on t\'ecnica}:
	arquitectura modular SPA + API,
	persistencia versionada con Flyway,
	seguridad JWT y soporte de integraciones cloud configurables.
	\item \textbf{Aportaci\'on de calidad}:
	incorporaci\'on de testing multinivel y mutaci\'on como criterio objetivo de robustez.
	\item \textbf{Aportaci\'on de operaci\'on}:
	canal CI/CD reproducible con validaciones continuas de calidad y seguridad.
\end{enumerate}

\section{Resultados y evidencia}

La evidencia disponible en el proyecto permite sostener de forma objetiva
la madurez alcanzada:

\begin{itemize}
	\item cobertura alta en frontend con umbrales exigentes en l\'ineas,
	funciones, sentencias y ramas;
	\item resultados favorables en mutaci\'on (frontend y backend),
	lo que indica que los tests detectan alteraciones de comportamiento;
	\item validaci\'on de flujos de usuario cr\'iticos mediante E2E;
	\item automatizaci\'on de verificaciones en pipelines de integraci\'on continua.
\end{itemize}

En conjunto, estos resultados reducen el riesgo de regresi\'on y
proporcionan una base fiable para mantenimiento y evoluci\'on del sistema.

Las decisiones de arquitectura, pruebas y seguridad se han contrastado
con documentaci\'on t\'ecnica oficial de las tecnolog\'ias empleadas y
recomendaciones de referencia en seguridad web
\cite{springbootdocs,reactdocs,flywaydocs,vitestdocs,playwrightdocs,owasptop10}.

\section{Evaluaci\'on del impacto por tipo de usuario}

El valor del proyecto se aprecia de forma diferente seg\'un el perfil de uso.
Por ello, se realiza una evaluaci\'on diferenciada:

\begin{itemize}
	\item \textbf{Profesorado}: disminuci\'on de tareas manuales repetitivas,
	mejor visibilidad del estado de entregas y consolidaci\'on del flujo de correcci\'on.
	\item \textbf{Estudiantado}: proceso de entrega m\'as guiado,
	mejor comprensi\'on del estado de su trabajo y acceso m\'as claro al feedback.
	\item \textbf{Administraci\'on}: mayor control sobre estructura acad\'emica y roles,
	con operaciones de gobierno centralizadas.
\end{itemize}

Desde esta perspectiva,
el sistema no solo aporta funcionalidad t\'ecnica,
sino tambi\'en una mejora organizativa en la operativa acad\'emica.

\section{Validaci\'on del enfoque de ingenier\'ia adoptado}

Una conclusi\'on relevante del trabajo es que el enfoque metodol\'ogico aplicado
(requisitos formalizados, arquitectura modular, pruebas en capas y CI/CD)
ha sido adecuado para un proyecto de este alcance.

Los principales indicadores de validez del enfoque son:

\begin{enumerate}
	\item coherencia entre objetivos, requisitos, implementaci\'on y verificaci\'on;
	\item estabilidad de las iteraciones sin desviaciones cr\'iticas de calendario;
	\item capacidad de incorporar mejoras sin ruptura estructural del sistema.
\end{enumerate}

Esto refuerza la idea de que,
en proyectos software de car\'acter acad\'emico-profesional,
la disciplina metodol\'ogica es un factor diferencial
tan importante como la selecci\'on tecnol\'ogica.

\section{Limitaciones del trabajo}

Como en todo proyecto con alcance temporal acotado,
se identifican limitaciones que delimitan el estado actual:

\begin{itemize}
	\item no todas las funcionalidades \emph{Should/Could} del backlog
	est\'an cerradas al 100\% en el corte del TFG;
	\item ciertas capacidades avanzadas (p. ej., autenticaci\'on federada completa,
	2FA integral u observabilidad avanzada) se mantienen como evoluci\'on posterior;
	\item la estrategia E2E predominante en el ciclo actual es mockeada,
	por lo que conviene reforzar escenarios con backend real en staging.
\end{itemize}

Estas limitaciones no invalidan el cumplimiento del objetivo principal,
pero orientan de forma clara la hoja de ruta futura.

\section{Lecciones aprendidas}

Durante la ejecuci\'on del proyecto se consolidan varios aprendizajes
de inter\'es para trabajos similares:

\begin{itemize}
	\item dise\~nar la trazabilidad desde el inicio (requisito-backlog-c\'odigo-prueba)
	simplifica justificaci\'on t\'ecnica y auditor\'ia;
	\item la calidad debe construirse iterativamente,
	no solo en la fase final de cierre;
	\item desacoplar configuraci\'on por entorno reduce fricci\'on en migraciones de infraestructura;
	\item integrar seguridad y permisos en la l\'ogica de dominio evita defectos estructurales costosos.
\end{itemize}

\section{Riesgos residuales y deuda t\'ecnica controlada}

En el estado actual se identifican riesgos residuales,
asumidos de forma expl\'icita y controlada:

\begin{itemize}
	\item cobertura funcional incompleta en funcionalidades \emph{Should/Could};
	\item necesidad de ampliar E2E contra backend real para ciertos escenarios;
	\item dependencia parcial de configuraciones externas en integraciones cloud.
\end{itemize}

La deuda t\'ecnica asociada se considera \textbf{controlada},
porque est\'a identificada,
priorizada y conectada con una hoja de ruta de evoluci\'on.

\section{Desarrollos futuros}

Se proponen l\'ineas de evoluci\'on priorizadas por impacto y viabilidad.

\subsection{Corto plazo}

\begin{itemize}
	\item cerrar PBIs pendientes de usabilidad y reporting docente;
	\item ampliar E2E con entorno semirreal y pruebas de regresi\'on ampliadas;
	\item reforzar auditor\'ia de permisos y evidencias de seguridad en preproducci\'on.
\end{itemize}

\subsection{Medio plazo}

\begin{itemize}
	\item autenticaci\'on federada institucional completa y doble factor transversal;
	\item observabilidad avanzada (m\'etricas de negocio, trazas distribuidas y alertado);
	\item funcionalidades anal\'iticas para seguimiento de actividad y carga de correcci\'on.
\end{itemize}

\subsection{Largo plazo}

\begin{itemize}
	\item incorporaci\'on de m\'odulos de apoyo docente
	(detecci\'on de patrones de entrega, recomendaci\'on de revisi\'on, etc.);
	\item estudio de interoperabilidad con LMS institucionales;
	\item escalado multi-instituci\'on con pol\'iticas avanzadas de gobierno de datos.
\end{itemize}

\section{Hoja de ruta priorizada de continuidad}

Para facilitar la continuidad del producto,
se sintetiza una priorizaci\'on por horizonte temporal y retorno esperado.

\begin{table}[ht]
\centering
\begin{tabular}{|P{2.5cm}|P{5.2cm}|P{3.8cm}|}
\hline
\textbf{Horizonte} & \textbf{Iniciativas principales} & \textbf{Impacto esperado} \\
\hline
Corto plazo & cierre de PBIs pendientes, ampliaci\'on de E2E y refuerzo de auditor\'ia de permisos & mejora inmediata de robustez operativa \\
\hline
Medio plazo & federaci\'on de identidad, 2FA y observabilidad avanzada & aumento de seguridad y gobierno t\'ecnico \\
\hline
Largo plazo & interoperabilidad institucional y m\'odulos avanzados de apoyo docente & escalabilidad funcional y adopci\'on ampliada \\
\hline
\end{tabular}
\caption{Hoja de ruta priorizada para evoluci\'on del sistema}
\label{tab:hoja-ruta-conclusiones}
\end{table}

\section{Reflexi\'on final}

El trabajo realizado demuestra que,
incluso en un marco temporal exigente,
es posible construir una soluci\'on acad\'emica con criterios de ingenier\'ia profesionales
si se mantiene una estrategia de trazabilidad,
priorizaci\'on de alcance y verificaci\'on continua.

La principal fortaleza del resultado no reside \'unicamente en las funcionalidades implementadas,
sino en la consistencia entre problema, decisiones adoptadas y evidencia de calidad.

\section{Cierre}

El TFG demuestra que es posible construir,
con un enfoque de ingenier\'ia disciplinado,
una soluci\'on acad\'emica de valor real para usuarios finales,
sin renunciar a criterios de calidad, seguridad y mantenibilidad.

La plataforma resultante ofrece una base s\'olida para continuidad,
adopci\'on progresiva y extensi\'on funcional,
y cumple con los objetivos acad\'emicos y t\'ecnicos definidos al inicio del proyecto.






